\section{Design}
\label{sec:design}

\begin{figure}[t]
\centering
\begin{tikzpicture}[
  font=\footnotesize,
  box/.style={draw=ink2, rounded corners=2pt, line width=0.5pt, align=center, inner sep=3pt, text width=2.55cm, minimum height=1.2cm, fill=white},
  llm/.style={box, fill=fillB},
  sone/.style={box, fill=fillA},
  store/.style={box, fill=fillC},
  flow/.style={-{Stealth[length=4.5pt]}, line width=0.6pt, draw=ink2},
  lab/.style={font=\scriptsize, text=ink2},
]
% main instance (left)
\node[box, text width=2.0cm, minimum height=0.8cm] (user) at (0,0) {user message};
\node[box, text width=2.0cm, minimum height=1.0cm] (main) at (0,-2.7) {main agent\\(answering model)};
\node[box, text width=2.0cm, minimum height=0.8cm] (reply) at (0,-4.8) {reply};
\draw[flow] (user) -- (main);
\draw[flow] (main) -- (reply);
% replica: reading
\node[llm] (plan) at (3.7,0) {\textbf{plan} (System 2)\\3 searches + a note;\\$\le$3 needs};
\node[store] (retrieve) at (7.0,0) {\textbf{retrieve}\\BM25 and vectors\\(RRF); pool of 48};
\node[sone] (screen) at (10.3,0) {\textbf{screen} (System 1)\\``should the reply\\use this record?''};
\node[sone] (judge) at (10.3,-2.7) {\textbf{judge} (System 1)\\needed? no longer\\current? needs? items?};
\node[box] (loop) at (7.0,-2.7) {\textbf{loop} ($\le$3 rounds)\\open needs: page,\\rewrite, or stop};
\node[box] (compose) at (3.7,-2.7) {\textbf{compose View}\\index items, then\\$\le$16 records / 12k};
% replica: memory
\node[store] (journal) at (7.0,2.2) {\textbf{journal Source}\\raw dated records};
\node[llm] (consolidate) at (10.3,2.2) {\textbf{consolidate}\\background, once\\per record};
\node[store] (index) at (13.3,2.2) {\textbf{index}\\timelines, values,\\instructions};
\draw[flow] (user) -- (plan);
\draw[flow] (plan) -- (retrieve);
\draw[flow] (retrieve) -- (screen);
\draw[flow] (screen) -- (judge);
\draw[flow] (judge) -- (loop);
\draw[flow] ([xshift=-4pt]loop.north) -- node[lab, left, align=right] {page /\\rewrite} ([xshift=-4pt]retrieve.south);
\draw[flow] (loop) -- node[lab, above] {stop} (compose);
\draw[flow, dashed] (journal) -- (retrieve);
\draw[flow, dashed] (journal) -- (consolidate);
\draw[flow, dashed] (consolidate) -- (index);
\draw[flow, dashed] (index.south) |- node[lab, pos=0.25, right, align=left] {items linked\\to the records} (judge.east);
\draw[flow] (compose) -- node[lab, above] {View} (main);
% frames
\begin{scope}[on background layer]
\node[draw=rule, dashed, rounded corners=4pt, fit=(user)(main)(reply), inner sep=6pt, label={[lab, font=\scriptsize\bfseries]above:{main DSH instance}}] {};
\node[draw=rule, dashed, rounded corners=4pt, fit=(plan)(retrieve)(screen)(judge)(loop)(compose)(journal)(consolidate)(index), inner sep=6pt, label={[lab, font=\scriptsize\bfseries]above:{replica DSH instance: \sys}}] {};
\end{scope}
% legend
\node[lab, anchor=north west, align=left] at (2.4,-3.75) {\fcolorbox{ink2}{fillB}{\rule{0pt}{4pt}\rule{4pt}{0pt}} System~2 (LLM) \quad \fcolorbox{ink2}{fillA}{\rule{0pt}{4pt}\rule{4pt}{0pt}} System~1 (\jev) \quad \fcolorbox{ink2}{fillC}{\rule{0pt}{4pt}\rule{4pt}{0pt}} storage and search \quad \fcolorbox{ink2}{white}{\rule{0pt}{4pt}\rule{4pt}{0pt}} rules \quad dashed: off the answer path};
\end{tikzpicture}
\caption{\sys{} at one user turn. System~2, an LLM, plans the searches; System~1, \jev, screens and judges the raw records they return and the index items those records link to; rules loop on needs that are still open and publish a View into the main instance, whose LLM answers. Off the answer path, System~2 consolidates each record once into an index that points back to the records.}
\label{fig:architecture}
\end{figure}

At every user turn \sys{} decides what the main agent should see before it replies (\cref{fig:architecture}).
It divides the work as \cref{sec:setting} does: an LLM plans the searches and names what the reply needs (System~2), \jev{} judges every record the searches return (System~1), rules with explicit budgets connect the two, and the System~2 work that no reply can wait for, consolidating the history into an index, runs in the background.
\Cref{sec:memory,sec:read,sec:rules} describe the memory, the two systems that read it and the rules; \cref{alg:turn} summarizes one turn.

\subsection{Memory: Raw Records and a Consolidated Index}
\label{sec:memory}

\paragraph{Records.}
Each session is written as it happened, as short \texttt{speaker: text} chunks, each headed by the session's number and date.
No LLM is on the write path of a record, and nothing about it is decided when it is written; the only computation is one embedding, so the write path is the same whatever the records are about.

\paragraph{Consolidation.}
Questions about a whole conversation, such as every mention of a topic, the current value of something that changed or an instruction given once, need evidence that the question's own searches rarely find.
Gathering it is System~2 work that no reply can wait for, so a background job reads the journal once, in order, and a small LLM folds it, batch by batch, into an \emph{index} of three kinds of items, each linked to the records it came from: events on topic timelines; dated histories of values that can change, marked where a value contradicts an earlier one; and standing instructions and preferences.
Items point to records and never replace them: the index is a view over the records, not their schema, and \sys{} reads the records with or without it.

\subsection{Reading with Two Systems}
\label{sec:read}

\paragraph{Plan (System~2).}
Planning is System~2 work: it must anticipate how the records that answer a question are worded and what the reply will have to say.
Two LLM calls read the last six messages in parallel.
The \emph{cue plan} writes three searches worded as a stored record would word them, plus one invented record that would answer the message~\citep{gao2023hyde}; with the message itself, a turn starts with five searches.
The \emph{need plan} lists at most three facts the reply needs, marked \texttt{NEED}, or \texttt{ALL} when the reply needs every instance of something (a count, a total, a list).
Neither call reads the memory, so their cost does not grow with it.

\paragraph{Retrieve.}
Each search ranks the journal's records by BM25 and by embedding similarity, fuses the two rankings by reciprocal rank~\citep{cormack2009rrf}, and returns a page of records.
The pages merge into a \emph{pool}: each search first gets an equal share of its own best results, and the rest go to the highest summed reciprocal rank.

\paragraph{Judge (System~1).}
Judging is System~1 work: each proposition concerns one record, has an explicit criterion and can be answered independently of the others.
\jev{} answers typed yes/no propositions and returns the probability of yes~\citep{jev2026}.
It \emph{screens} each pooled record (should the assistant's next reply use this item?), then \emph{judges} the records that could enter the View, beside each other and the last six messages, as \emph{needed} or \emph{no longer current} (a later message or record explicitly changes, cancels or completes it), and fills a \emph{need table}: does this record give what this need asks for?
The judged records light up the index items linked to them, and \jev{} judges which of these topics and values, and of the dozen nearest to the question, the reply needs.
Counting, arithmetic, comparing dates and multi-hop indirection are System~2 work, which \jev{} handles poorly~\citep{jevdocs2026}; they are left to the planner and the answering model.

\subsection{Rules and Budgets}
\label{sec:rules}

The rules connect the two systems.
They read only what System~1 reliably gives, its yes/no line at 0.5 and its order, together with events observed in the previous step, and they call on System~2 only when System~1 reports trouble: a need that no judged record gives prompts the planner for a new search, much as System~2 is mobilized when System~1 runs into difficulty~\citep{kahneman2011thinking}.
Every other number is a budget with a meaning: three needs per turn, pages of 20 records or 12,000 characters, a pool of 48 records, three rounds, and a View of 16 records and 12,000 characters.

\paragraph{The loop.}
Each open need gets one action per round.
A single need is met once a record gives it.
A need for every instance, labelled \texttt{ALL} or observed because two records give it, pages each search that found a new yes until none appears.
A need that no record gives yet follows its \emph{lead}, the record the need table ranks highest: the loop pages the search that ranked the lead highest, then asks the planner for one new search, then gives the need up.

\paragraph{The View.}
The View leads with the index, each block within a fixed budget: standing instructions, a one-line directory of topics, and the timelines and value histories \jev{} judged needed.
Records follow: those judged needed, then, when the planner named a need, the rest in \jev's order up to the View's budget, trimmed from the end to make room for the index.
Records judged no longer current are left out.

\begin{algorithm}[t]
\caption{One user turn of \sys.}
\label{alg:turn}
\small
\begin{algorithmic}[1]
\Require the last six messages $M$, the journal $J$ and its index $X$
\State $(Q, N) \gets \textsc{Plan}(M)$ \Comment{System~2: five searches, at most three needs}
\State $P \gets \textsc{Pool}(\{\textsc{Search}(J, q) : q \in Q\})$ \Comment{BM25 and vectors, fused; at most 48 records}
\For{round $= 1, 2, 3$}
  \State $\textsc{Screen}(P)$;\; $\textsc{Judge}(P, M, N)$ \Comment{System~1: use it? needed? no longer current? which need?}
  \State $A \gets \textsc{Act}(N, P)$ \Comment{rules: page a search, ask System~2 for a new one, or give up}
  \If{$A = \emptyset$} \State \textbf{break} \EndIf
  \State $P \gets P \cup \textsc{Run}(A)$
\EndFor
\State $I \gets \textsc{JudgeItems}(X, P, M)$ \Comment{index items linked to judged records, and the nearest}
\State \Return $\textsc{Compose}(I, P)$ \Comment{index first, then records, within the View's budgets}
\end{algorithmic}
\end{algorithm}
